% GENERATED FILE. Edit canonical lessons, then run npm run generate:books.
% canonical-source-hash: fnv1a64:ad21d208927b49c7
% canonical-lessons: TA-C199-nampu, TA-C199-ottukkol, TA-C199-maru, TA-C199-vivati, TA-C199-oppitu

\chapter{To believe, To agree, To refuse, To discuss, To compare}
\label{ch:ta-to-believe-to-agree-to-refuse-to-discuss-to-compare}
\hlchaptermodality{199}

\begin{chapteropening}
\textbf{By the end of this chapter:} \emph{I can say to believe, to agree, to refuse, to discuss and to compare.}

Complete the last lesson of chapter 199: I can say to believe, to agree, to refuse, to discuss and to compare.
\end{chapteropening}

\section[\texorpdfstring{\ta{நம்பு} (nampu)}{நம்பு (nampu)}]{\ta{நம்பு} (nampu) — to believe}

\label{lesson:TA-C199-nampu}



\begin{quote}
Before the new one: say the Tamil for to quarrel, then the Tamil for to tease.
\end{quote}

\subsection*{You'll want to know: \ta{நம்பு}}
\textbf{\ta{நம்பு}} — \emph{nampu} — \textquotedblleft{}to believe\textquotedblright{}.

\begin{grammarlens}[title={What you've built}]
One more word for this chapter.
\end{grammarlens}

\subsection*{Your turn}
\begin{itemize}
\raggedright
  \item \emph{Say it:} \emph{nampu}
  \item \emph{Say it:} \emph{nampu}, once more
  \item \emph{Say it:} say \emph{nampu}
  \item \emph{Recall:} say the Tamil for to quarrel, then the Tamil for to tease, then say \emph{nampu} again
\end{itemize}

\begin{tcolorbox}[breakable,colback=teal!4,colframe=teal!35!black,title={Before you move on}]
What does \ta{நம்பு} mean? (\textquotedblleft{}To believe\textquotedblright{}.) Say it once more.
\end{tcolorbox}
\section[\texorpdfstring{\ta{ஒத்துக்கொள்} (ottukkoḷ)}{ஒத்துக்கொள் (ottukkoḷ)}]{\ta{ஒத்துக்கொள்} (ottukkoḷ) — to agree, to admit}

\label{lesson:TA-C199-ottukkol}



\begin{quote}
Before the new one: say the Tamil for to tease, then the Tamil for to believe.
\end{quote}

\subsection*{You'll want to know: \ta{ஒத்துக்கொள்}}
\textbf{\ta{ஒத்துக்கொள்}} — \emph{ottukkoḷ} — \textquotedblleft{}to agree, to admit\textquotedblright{}.

\begin{grammarlens}[title={What you've built}]
One more word for this chapter.
\end{grammarlens}

\subsection*{Your turn}
\begin{itemize}
\raggedright
  \item \emph{Say it:} \emph{ottukkoḷ}
  \item \emph{Say it:} \emph{ottukkoḷ}, once more
  \item \emph{Say it:} say \emph{ottukkoḷ}
  \item \emph{Recall:} say the Tamil for to tease, then the Tamil for to believe, then say \emph{ottukkoḷ} again
\end{itemize}

\begin{tcolorbox}[breakable,colback=teal!4,colframe=teal!35!black,title={Before you move on}]
What does \ta{ஒத்துக்கொள்} mean? (\textquotedblleft{}To agree, to admit\textquotedblright{}.) Say it once more.
\end{tcolorbox}
\section[\texorpdfstring{\ta{மறு} (maṟu)}{மறு (maṟu)}]{\ta{மறு} (maṟu) — to refuse, to deny}

\label{lesson:TA-C199-maru}



\begin{quote}
Before the new one: say the Tamil for to believe, then the Tamil for to agree.
\end{quote}

\subsection*{You'll want to know: \ta{மறு}}
\textbf{\ta{மறு}} — \emph{maṟu} — \textquotedblleft{}to refuse, to deny\textquotedblright{}.

\begin{grammarlens}[title={What you've built}]
One more word for this chapter.
\end{grammarlens}

\subsection*{Your turn}
\begin{itemize}
\raggedright
  \item \emph{Say it:} \emph{maṟu}
  \item \emph{Say it:} \emph{maṟu}, once more
  \item \emph{Say it:} say \emph{maṟu}
  \item \emph{Recall:} say the Tamil for to believe, then the Tamil for to agree, then say \emph{maṟu} again
\end{itemize}

\begin{tcolorbox}[breakable,colback=teal!4,colframe=teal!35!black,title={Before you move on}]
What does \ta{மறு} mean? (\textquotedblleft{}To refuse, to deny\textquotedblright{}.) Say it once more.
\end{tcolorbox}
\section[\texorpdfstring{\ta{விவாதி} (vivāti)}{விவாதி (vivāti)}]{\ta{விவாதி} (vivāti) — to discuss, to debate}

\label{lesson:TA-C199-vivati}



\begin{quote}
Before the new one: say the Tamil for to agree, then the Tamil for to refuse.
\end{quote}

\subsection*{You'll want to know: \ta{விவாதி}}
\textbf{\ta{விவாதி}} — \emph{vivāti} — \textquotedblleft{}to discuss, to debate\textquotedblright{}.

\begin{grammarlens}[title={What you've built}]
One more word for this chapter.
\end{grammarlens}

\subsection*{Your turn}
\begin{itemize}
\raggedright
  \item \emph{Say it:} \emph{vivāti}
  \item \emph{Say it:} \emph{vivāti}, once more
  \item \emph{Say it:} say \emph{vivāti}
  \item \emph{Recall:} say the Tamil for to agree, then the Tamil for to refuse, then say \emph{vivāti} again
\end{itemize}

\begin{tcolorbox}[breakable,colback=teal!4,colframe=teal!35!black,title={Before you move on}]
What does \ta{விவாதி} mean? (\textquotedblleft{}To discuss, to debate\textquotedblright{}.) Say it once more.
\end{tcolorbox}
\section[\texorpdfstring{\ta{ஒப்பிடு} (oppiṭu)}{ஒப்பிடு (oppiṭu)}]{\ta{ஒப்பிடு} (oppiṭu) — to compare}

\label{lesson:TA-C199-oppitu}



\begin{quote}
Before the new one: say the Tamil for to refuse, then the Tamil for to discuss.
\end{quote}

\subsection*{You'll want to know: \ta{ஒப்பிடு}}
\textbf{\ta{ஒப்பிடு}} — \emph{oppiṭu} — \textquotedblleft{}to compare\textquotedblright{}.

\begin{grammarlens}[title={What you've built}]
One more word for this chapter.
\end{grammarlens}

\subsection*{Your turn}
\begin{itemize}
\raggedright
  \item \emph{Say it:} \emph{oppiṭu}
  \item \emph{Say it:} \emph{oppiṭu}, once more
  \item \emph{Say it:} say \emph{oppiṭu}
  \item \emph{Recall:} say the Tamil for to refuse, then the Tamil for to discuss, then say \emph{oppiṭu} again
\end{itemize}

\begin{tcolorbox}[breakable,colback=teal!4,colframe=teal!35!black,title={Before you move on}]
What does \ta{ஒப்பிடு} mean? (\textquotedblleft{}To compare\textquotedblright{}.) Say it once more.
\end{tcolorbox}
